\section{Introduction}


The emergence of large language models (LLMs) as reasoning engines has created demand for robust frameworks that can orchestrate multiple specialized agents across complex workflows. While existing agent frameworks provide basic tool-calling capabilities, they often lack: (1) type-safe context propagation across agent boundaries, (2) memory systems that support both short-term and long-term recall, and (3) mechanisms for intelligent task routing in multi-agent scenarios.

The Hanzo Agent SDK addresses these limitations through a principled architecture that treats agents as first-class, composable units. Each agent is defined as a generic dataclass parameterized by a context type, enabling type-safe state propagation throughout the execution graph. The framework supports dynamic instruction generation, input/output guardrails, and hierarchical handoffs between specialized agents.

\paragraph{Contributions.} This paper makes the following contributions:
\begin{itemize}
    \item A generic \texttt{Agent[TContext]} abstraction with support for dynamic instructions, model-specific settings, and composable guardrails.
    \item A memory system with three storage tiers (conversation, reflection, fact) supporting semantic retrieval and automatic compression.
    \item An \texttt{AgentNetwork} for multi-agent orchestration with semantic routing, shared state, and parallel execution.
    \item Comprehensive benchmarks demonstrating 34\% improvement in task completion and 2.1x reduction in token usage through intelligent routing.
\end{itemize}

